\documentclass[11pt,a4paper]{article}
\usepackage[utf8]{inputenc}
\usepackage[T1]{fontenc}
\usepackage[catalan]{babel}
\usepackage{lmodern}
\usepackage[sfdefault,lining]{FiraSans}
\usepackage[a4paper,top=2.3cm,bottom=2.2cm,left=1.8cm,right=1.8cm,headheight=1.3cm,headsep=0.4cm,footskip=1cm]{geometry}
\usepackage{amsmath,amssymb}
\usepackage{icomma}
\usepackage{graphicx}
\usepackage[table]{xcolor}
\usepackage{tikz}
\usepackage{fontawesome5}
\usepackage{ccicons}
\usepackage{fancyhdr}
\usepackage{lastpage}
\usepackage{tasks}
\usepackage{needspace}
\usepackage[most]{tcolorbox}
\usepackage{array}
\usepackage[hidelinks,pdftitle={Fitxa d'exercicis de la unitat 8: Semblança. Teorema de Tales. 2n d'ESO},pdfsubject={Matemàtiques 2n ESO},pdfkeywords={semblança, figures semblants, raó de semblança, escala, plànols, mapes, maquetes, teorema de Tales, posició de Tales, criteris de semblança, ombres}]{hyperref}

\definecolor{navy}{HTML}{1B2748}
\definecolor{accent}{HTML}{E8603A}
\definecolor{teal}{HTML}{1F9E90}
\definecolor{soft}{HTML}{EEF1F8}
\definecolor{muted}{HTML}{5D6685}

\setlength{\parindent}{0pt}
\setlength{\parskip}{3pt}
\settasks{label=\alph*), label-format=\bfseries\color{navy}, label-width=1.3em, item-indent=1.8em, column-sep=0.6em, after-item-skip=5pt, before-skip=3pt, after-skip=4pt}

% ---------- capçalera i peu ----------
\newcommand{\llicencia}{\href{https://creativecommons.org/licenses/by-nc-sa/4.0/deed.ca}{\ccbyncsa}\ \ Llicència Creative Commons Reconeixement-NoComercial-CompartirIgual 4.0 Internacional (CC BY-NC-SA 4.0)}
\pagestyle{fancy}
\fancyhf{}
\fancyhead[L]{\raisebox{-0.35\height}{\includegraphics[height=0.95cm]{logo-gva.png}}}
\fancyhead[C]{\small\color{muted}Fitxa d'exercicis · Unitat 8: Semblança. Teorema de Tales · 2n d'ESO}
\fancyhead[R]{\raisebox{-0.35\height}{\includegraphics[height=0.95cm]{logo-vermellar.png}}}
\fancyfoot[L]{\scriptsize\color{muted}\llicencia}
\fancyfoot[R]{\small\color{muted}\thepage}
\renewcommand{\headrulewidth}{0.4pt}
\renewcommand{\headrule}{\color{muted!40}\hrule height \headrulewidth}
\fancypagestyle{primera}{\fancyhf{}\renewcommand{\headrulewidth}{0pt}
  \fancyfoot[L]{\scriptsize\color{muted}\llicencia}
  \fancyfoot[R]{\small\color{muted}\thepage}}

% ---------- nivells i exercicis ----------
\newcommand{\estrelles}[1]{{\color{accent}\ifcase#1\or\faStar\faStar[regular]\faStar[regular]\or\faStar\faStar\faStar[regular]\or\faStar\faStar\faStar\fi}}
\newcounter{exer}
\newcommand{\ex}[1]{\par\needspace{4\baselineskip}\medskip\stepcounter{exer}\noindent{\color{navy}\bfseries\theexer.}\ \ {\small\estrelles{#1}}\ \ }

\newcommand{\seccio}[2]{\par\needspace{10\baselineskip}\bigskip
  \begin{tcolorbox}[colback=navy,colframe=navy,coltext=white,arc=2mm,boxrule=0pt,left=3mm,right=3mm,top=1.2mm,bottom=1.2mm]
  {\large\bfseries #1\quad #2}
  \end{tcolorbox}\vspace{-2pt}}

\newcolumntype{C}[1]{>{\centering\arraybackslash}p{#1}}
\usepackage{pifont}
% quadre de recordatori (només a la fitxa d'exercicis)
\newcommand{\recorda}[1]{\par\smallskip\begin{tcolorbox}[colback=accent!7,colframe=accent!55,arc=2mm,boxrule=0.5pt,left=2.5mm,right=2.5mm,top=1mm,bottom=1mm]\small{\color{accent}\faLightbulb[regular]}\ \textbf{Recorda:}\ #1\end{tcolorbox}}
\newcommand{\res}[1]{#1}\newcommand{\rest}[1]{#1}

\def\fmneed{14}
% =====================================================================
% Capçalera FentMates de les fitxes: un color per a cada curs.
% Els colors són els de la pàgina inicial (LEVELS a assets/js/materials.js);
% el logo ve de _eines/logo/pdf/ (python3 _eines/logo/genera.py).
%
% Ús, al final del preàmbul de la unitat:
%   % =====================================================================
% Capçalera FentMates de les fitxes: un color per a cada curs.
% Els colors són els de la pàgina inicial (LEVELS a assets/js/materials.js);
% el logo ve de _eines/logo/pdf/ (python3 _eines/logo/genera.py).
%
% Ús, al final del preàmbul de la unitat:
%   % =====================================================================
% Capçalera FentMates de les fitxes: un color per a cada curs.
% Els colors són els de la pàgina inicial (LEVELS a assets/js/materials.js);
% el logo ve de _eines/logo/pdf/ (python3 _eines/logo/genera.py).
%
% Ús, al final del preàmbul de la unitat:
%   \input{../comu/fentmates.tex}
%   \fmcurs{1eso}
% i, al començament del document:
%   \capcalera{Matemàtiques}{1r ESO}{Fitxa d'exercicis · Unitat 2}{Divisibilitat}
% Opcional abans de \input: \def\fmneed{14} (espai mínim abans de cada secció).
% =====================================================================
\makeatletter
\def\fm@logodir{../../logo/pdf}
\providecommand{\fmneed}{10}
\def\fm@cursnom{}

% clau del curs, color i color del text damunt del color
\newcommand{\fm@curs}[3]{\@namedef{fm@color@#1}{#2}\@namedef{fm@text@#1}{#3}}
\fm@curs{1eso}{E6B84A}{navy}
\fm@curs{2eso}{3B82F6}{white}
\fm@curs{3eso}{FF7A50}{white}
\fm@curs{4eso-a}{A78BFA}{white}
\fm@curs{4eso-b}{6D28D9}{white}
\fm@curs{1batx-mi}{9B1C1C}{white}
\fm@curs{1batx-ccss}{EF6B6B}{white}
\fm@curs{2batx-mii}{1E6B47}{white}
\fm@curs{2batx-ccss}{2BB3A3}{white}

\newcommand{\fmcurs}[1]{%
  \@ifundefined{fm@color@#1}{\PackageError{fentmates}{Curs desconegut: #1}{}}{}%
  \edef\fm@tmp{\noexpand\definecolor{curs}{HTML}{\@nameuse{fm@color@#1}}%
               \noexpand\colorlet{curstext}{\@nameuse{fm@text@#1}}}\fm@tmp
  \colorlet{cursfosc}{curs!72!black}%
  \def\fm@logo{\fm@logodir/vertical-#1.pdf}%
  \def\fm@marca{\fm@logodir/marca-#1.pdf}}

% ---------- capçalera de la primera pàgina ----------
% #1 matèria (petit), #2 curs, #3 tipus de document i unitat, #4 títol
\newcommand{\capcalera}[4]{%
  \gdef\fm@cursnom{#2}%
  \par\noindent
  \begin{tcolorbox}[colback=curs,colframe=curs,boxrule=0pt,arc=2.2mm,width=3.5cm,height=2.15cm,
      valign=center,left=3mm,right=2mm,top=1mm,bottom=1mm,nobeforeafter,box align=center]
    \color{curstext}{\footnotesize #1}\par\vspace{1pt}{\fontsize{19}{21}\selectfont\bfseries #2\par}%
  \end{tcolorbox}\hspace{5mm}%
  \begin{minipage}[c]{\dimexpr\linewidth-3.5cm-5mm-2.9cm\relax}
    {\color{cursfosc}\small\bfseries\MakeUppercase{#3}}\par\vspace{2pt}
    {\color{navy}\fontsize{20}{23}\selectfont\bfseries #4\par}
  \end{minipage}\hfill
  \begin{minipage}[c]{2.6cm}\raggedleft\includegraphics[width=2.6cm]{\fm@logo}\end{minipage}\par
  \vspace{2pt}}

% ---------- barres de secció del color del curs ----------
\renewcommand{\seccio}[2]{\par\needspace{\fmneed\baselineskip}\bigskip
  \begin{tcolorbox}[colback=curs,colframe=curs,coltext=curstext,arc=2mm,boxrule=0pt,left=3mm,right=3mm,top=1.2mm,bottom=1.2mm]
  {\large\bfseries #1\quad #2}
  \end{tcolorbox}\vspace{-2pt}}

% ---------- estrelles de dificultat del color del curs (a Batxillerat no se'n posen) ----------
\@ifundefined{estrelles}{}{%
  \renewcommand{\estrelles}[1]{{\color{curs}\ifcase#1\or\faStar\faStar[regular]\faStar[regular]\or\faStar\faStar\faStar[regular]\or\faStar\faStar\faStar\fi}}}

% ---------- capçalera i peu de les pàgines ----------
% primera pàgina: logos institucionals al peu
\fancypagestyle{primera}{\fancyhf{}\renewcommand{\headrulewidth}{0pt}
  \fancyfoot[L]{\raisebox{-0.32\height}{\includegraphics[height=0.78cm]{logo-gva.png}}\hspace{5mm}%
    \raisebox{-0.32\height}{\includegraphics[height=0.78cm]{logo-vermellar.png}}}
  \fancyfoot[R]{\scriptsize\color{muted}\href{https://creativecommons.org/licenses/by-nc-sa/4.0/deed.ca}{\ccbyncsa}\ \ CC BY-NC-SA 4.0 · fentmates.org\qquad\small\thepage}}
% la resta: el símbol a l'esquerra i el curs a la dreta, amb el títol al centre com abans
\fancyhead[L]{\raisebox{-0.3\height}{\includegraphics[height=0.62cm]{\fm@marca}}}
\fancyhead[R]{\small\bfseries\color{cursfosc}\fm@cursnom}
\renewcommand{\headrule}{\color{curs!45}\hrule height \headrulewidth}
\makeatother

%   \fmcurs{1eso}
% i, al començament del document:
%   \capcalera{Matemàtiques}{1r ESO}{Fitxa d'exercicis · Unitat 2}{Divisibilitat}
% Opcional abans de \input: \def\fmneed{14} (espai mínim abans de cada secció).
% =====================================================================
\makeatletter
\def\fm@logodir{../../logo/pdf}
\providecommand{\fmneed}{10}
\def\fm@cursnom{}

% clau del curs, color i color del text damunt del color
\newcommand{\fm@curs}[3]{\@namedef{fm@color@#1}{#2}\@namedef{fm@text@#1}{#3}}
\fm@curs{1eso}{E6B84A}{navy}
\fm@curs{2eso}{3B82F6}{white}
\fm@curs{3eso}{FF7A50}{white}
\fm@curs{4eso-a}{A78BFA}{white}
\fm@curs{4eso-b}{6D28D9}{white}
\fm@curs{1batx-mi}{9B1C1C}{white}
\fm@curs{1batx-ccss}{EF6B6B}{white}
\fm@curs{2batx-mii}{1E6B47}{white}
\fm@curs{2batx-ccss}{2BB3A3}{white}

\newcommand{\fmcurs}[1]{%
  \@ifundefined{fm@color@#1}{\PackageError{fentmates}{Curs desconegut: #1}{}}{}%
  \edef\fm@tmp{\noexpand\definecolor{curs}{HTML}{\@nameuse{fm@color@#1}}%
               \noexpand\colorlet{curstext}{\@nameuse{fm@text@#1}}}\fm@tmp
  \colorlet{cursfosc}{curs!72!black}%
  \def\fm@logo{\fm@logodir/vertical-#1.pdf}%
  \def\fm@marca{\fm@logodir/marca-#1.pdf}}

% ---------- capçalera de la primera pàgina ----------
% #1 matèria (petit), #2 curs, #3 tipus de document i unitat, #4 títol
\newcommand{\capcalera}[4]{%
  \gdef\fm@cursnom{#2}%
  \par\noindent
  \begin{tcolorbox}[colback=curs,colframe=curs,boxrule=0pt,arc=2.2mm,width=3.5cm,height=2.15cm,
      valign=center,left=3mm,right=2mm,top=1mm,bottom=1mm,nobeforeafter,box align=center]
    \color{curstext}{\footnotesize #1}\par\vspace{1pt}{\fontsize{19}{21}\selectfont\bfseries #2\par}%
  \end{tcolorbox}\hspace{5mm}%
  \begin{minipage}[c]{\dimexpr\linewidth-3.5cm-5mm-2.9cm\relax}
    {\color{cursfosc}\small\bfseries\MakeUppercase{#3}}\par\vspace{2pt}
    {\color{navy}\fontsize{20}{23}\selectfont\bfseries #4\par}
  \end{minipage}\hfill
  \begin{minipage}[c]{2.6cm}\raggedleft\includegraphics[width=2.6cm]{\fm@logo}\end{minipage}\par
  \vspace{2pt}}

% ---------- barres de secció del color del curs ----------
\renewcommand{\seccio}[2]{\par\needspace{\fmneed\baselineskip}\bigskip
  \begin{tcolorbox}[colback=curs,colframe=curs,coltext=curstext,arc=2mm,boxrule=0pt,left=3mm,right=3mm,top=1.2mm,bottom=1.2mm]
  {\large\bfseries #1\quad #2}
  \end{tcolorbox}\vspace{-2pt}}

% ---------- estrelles de dificultat del color del curs (a Batxillerat no se'n posen) ----------
\@ifundefined{estrelles}{}{%
  \renewcommand{\estrelles}[1]{{\color{curs}\ifcase#1\or\faStar\faStar[regular]\faStar[regular]\or\faStar\faStar\faStar[regular]\or\faStar\faStar\faStar\fi}}}

% ---------- capçalera i peu de les pàgines ----------
% primera pàgina: logos institucionals al peu
\fancypagestyle{primera}{\fancyhf{}\renewcommand{\headrulewidth}{0pt}
  \fancyfoot[L]{\raisebox{-0.32\height}{\includegraphics[height=0.78cm]{logo-gva.png}}\hspace{5mm}%
    \raisebox{-0.32\height}{\includegraphics[height=0.78cm]{logo-vermellar.png}}}
  \fancyfoot[R]{\scriptsize\color{muted}\href{https://creativecommons.org/licenses/by-nc-sa/4.0/deed.ca}{\ccbyncsa}\ \ CC BY-NC-SA 4.0 · fentmates.org\qquad\small\thepage}}
% la resta: el símbol a l'esquerra i el curs a la dreta, amb el títol al centre com abans
\fancyhead[L]{\raisebox{-0.3\height}{\includegraphics[height=0.62cm]{\fm@marca}}}
\fancyhead[R]{\small\bfseries\color{cursfosc}\fm@cursnom}
\renewcommand{\headrule}{\color{curs!45}\hrule height \headrulewidth}
\makeatother

%   \fmcurs{1eso}
% i, al començament del document:
%   \capcalera{Matemàtiques}{1r ESO}{Fitxa d'exercicis · Unitat 2}{Divisibilitat}
% Opcional abans de \input: \def\fmneed{14} (espai mínim abans de cada secció).
% =====================================================================
\makeatletter
\def\fm@logodir{../../logo/pdf}
\providecommand{\fmneed}{10}
\def\fm@cursnom{}

% clau del curs, color i color del text damunt del color
\newcommand{\fm@curs}[3]{\@namedef{fm@color@#1}{#2}\@namedef{fm@text@#1}{#3}}
\fm@curs{1eso}{E6B84A}{navy}
\fm@curs{2eso}{3B82F6}{white}
\fm@curs{3eso}{FF7A50}{white}
\fm@curs{4eso-a}{A78BFA}{white}
\fm@curs{4eso-b}{6D28D9}{white}
\fm@curs{1batx-mi}{9B1C1C}{white}
\fm@curs{1batx-ccss}{EF6B6B}{white}
\fm@curs{2batx-mii}{1E6B47}{white}
\fm@curs{2batx-ccss}{2BB3A3}{white}

\newcommand{\fmcurs}[1]{%
  \@ifundefined{fm@color@#1}{\PackageError{fentmates}{Curs desconegut: #1}{}}{}%
  \edef\fm@tmp{\noexpand\definecolor{curs}{HTML}{\@nameuse{fm@color@#1}}%
               \noexpand\colorlet{curstext}{\@nameuse{fm@text@#1}}}\fm@tmp
  \colorlet{cursfosc}{curs!72!black}%
  \def\fm@logo{\fm@logodir/vertical-#1.pdf}%
  \def\fm@marca{\fm@logodir/marca-#1.pdf}}

% ---------- capçalera de la primera pàgina ----------
% #1 matèria (petit), #2 curs, #3 tipus de document i unitat, #4 títol
\newcommand{\capcalera}[4]{%
  \gdef\fm@cursnom{#2}%
  \par\noindent
  \begin{tcolorbox}[colback=curs,colframe=curs,boxrule=0pt,arc=2.2mm,width=3.5cm,height=2.15cm,
      valign=center,left=3mm,right=2mm,top=1mm,bottom=1mm,nobeforeafter,box align=center]
    \color{curstext}{\footnotesize #1}\par\vspace{1pt}{\fontsize{19}{21}\selectfont\bfseries #2\par}%
  \end{tcolorbox}\hspace{5mm}%
  \begin{minipage}[c]{\dimexpr\linewidth-3.5cm-5mm-2.9cm\relax}
    {\color{cursfosc}\small\bfseries\MakeUppercase{#3}}\par\vspace{2pt}
    {\color{navy}\fontsize{20}{23}\selectfont\bfseries #4\par}
  \end{minipage}\hfill
  \begin{minipage}[c]{2.6cm}\raggedleft\includegraphics[width=2.6cm]{\fm@logo}\end{minipage}\par
  \vspace{2pt}}

% ---------- barres de secció del color del curs ----------
\renewcommand{\seccio}[2]{\par\needspace{\fmneed\baselineskip}\bigskip
  \begin{tcolorbox}[colback=curs,colframe=curs,coltext=curstext,arc=2mm,boxrule=0pt,left=3mm,right=3mm,top=1.2mm,bottom=1.2mm]
  {\large\bfseries #1\quad #2}
  \end{tcolorbox}\vspace{-2pt}}

% ---------- estrelles de dificultat del color del curs (a Batxillerat no se'n posen) ----------
\@ifundefined{estrelles}{}{%
  \renewcommand{\estrelles}[1]{{\color{curs}\ifcase#1\or\faStar\faStar[regular]\faStar[regular]\or\faStar\faStar\faStar[regular]\or\faStar\faStar\faStar\fi}}}

% ---------- capçalera i peu de les pàgines ----------
% primera pàgina: logos institucionals al peu
\fancypagestyle{primera}{\fancyhf{}\renewcommand{\headrulewidth}{0pt}
  \fancyfoot[L]{\raisebox{-0.32\height}{\includegraphics[height=0.78cm]{logo-gva.png}}\hspace{5mm}%
    \raisebox{-0.32\height}{\includegraphics[height=0.78cm]{logo-vermellar.png}}}
  \fancyfoot[R]{\scriptsize\color{muted}\href{https://creativecommons.org/licenses/by-nc-sa/4.0/deed.ca}{\ccbyncsa}\ \ CC BY-NC-SA 4.0 · fentmates.org\qquad\small\thepage}}
% la resta: el símbol a l'esquerra i el curs a la dreta, amb el títol al centre com abans
\fancyhead[L]{\raisebox{-0.3\height}{\includegraphics[height=0.62cm]{\fm@marca}}}
\fancyhead[R]{\small\bfseries\color{cursfosc}\fm@cursnom}
\renewcommand{\headrule}{\color{curs!45}\hrule height \headrulewidth}
\makeatother

\fmcurs{2eso}

\begin{document}
\thispagestyle{primera}

\capcalera{Matemàtiques}{2n ESO}{Fitxa d'exercicis · Unitat 8}{Semblança. Teorema de Tales}
\vspace{10pt}
Nom i cognoms: \rule{8.4cm}{0.4pt}\hfill Grup: \rule{1.6cm}{0.4pt}\hfill Data: \rule{2.2cm}{0.4pt}
\vspace{8pt}

\begin{tcolorbox}[colback=soft,colframe=soft,arc=2mm,left=3mm,right=3mm,top=1.5mm,bottom=1.5mm]
\small Exercicis graduats per nivells de dificultat:\quad \estrelles{1}\ bàsic \qquad \estrelles{2}\ mitjà \qquad \estrelles{3}\ avançat.\\[2pt] Resol-los al quadern mostrant tots els passos i fes un dibuix de cada situació. Escriu la proporció abans de calcular. Les seccions segueixen l'ordre dels apunts de la unitat (8.1--8.5).
\end{tcolorbox}

% =====================================================
\seccio{1}{Figures semblants}

\ex{1} Digues si són sempre semblants (\textit{sí} o \textit{no}) i explica per què:\begin{tasks}(2)
  \task Dos quadrats qualssevol.
  \task Dos rectangles qualssevol.
  \task Dos triangles equilàters.
  \task Un quadrat i un rombe de 4 cm de costat.
\end{tasks}\recorda{Dues figures són \textbf{semblants} si tenen els angles iguals i els costats homòlegs proporcionals. La \textbf{raó de semblança} és $k=\frac{\text{costat nou}}{\text{costat original}}$: si $k>1$ és una ampliació i si $k<1$, una reducció.}

\ex{1} Calcula la raó de semblança i digues si és una ampliació o una reducció:\begin{tasks}(1)
  \task D'un rectangle de 4 cm × 6 cm a un altre de 10 cm × 15 cm.
  \task D'un triangle de costats 3, 4 i 5 cm a un altre de costats 1,5, 2 i 2,5 cm.
\end{tasks}

\ex{1} Un triangle té costats de 5, 7 i 9 cm. Calcula els costats d'un triangle semblant: \ a) amb raó $k=3$; \ b) amb raó $k=0{,}4$.

\ex{1} Un rectangle de 6 cm × 9 cm és semblant a un altre que té el costat menut de 4 cm. Calcula la raó de semblança i el costat gran del segon rectangle.

\ex{2} Les dues figures són semblants. Calcula la raó de semblança i els costats $x$ i $y$.\par\smallskip\begin{center}\begin{tikzpicture}[x=0.42cm,y=0.42cm,baseline=(current bounding box.center),line join=round]
\draw[fill=teal!12,draw=teal,thick] (0.000,0.000) -- (4.000,0.000) -- (4.000,1.000) -- (1.500,1.000) -- (1.500,3.000) -- (0.000,3.000) -- cycle;
\node[font=\small,inner sep=2pt,below] at (2.000,0.000) {$4$};
\node[font=\small,inner sep=2pt,right] at (4.000,0.500) {$1$};
\node[font=\small,inner sep=2pt,above] at (2.750,1.000) {$2{,}5$};
\node[font=\small,inner sep=2pt,right] at (1.500,2.000) {$2$};
\node[font=\small,inner sep=2pt,above] at (0.750,3.000) {$1{,}5$};
\node[font=\small,inner sep=2pt,left] at (0.000,1.500) {$3$};
\draw[fill=accent!22,draw=accent,thick] (6.000,0.000) -- (12.000,0.000) -- (12.000,1.500) -- (8.250,1.500) -- (8.250,4.500) -- (6.000,4.500) -- cycle;
\node[font=\small,inner sep=2pt,below] at (9.000,0.000) {$6$};
\node[font=\small,inner sep=2pt,right] at (12.000,0.750) {$x$};
\node[font=\small,inner sep=2pt,right] at (8.250,3.000) {$y$};
\node[font=\small,inner sep=2pt,left] at (6.000,2.250) {$4{,}5$};
\end{tikzpicture}\end{center}

\ex{2} Són semblants aquests rectangles? Si ho són, calcula la raó de semblança.\begin{tasks}(2)
  \task 4 cm × 10 cm \ i \ 6 cm × 15 cm
  \task 5 cm × 8 cm \ i \ 10 cm × 18 cm
\end{tasks}

\ex{2} Un polígon té 24 cm de perímetre i 32 cm² d'àrea. Calcula el perímetre i l'àrea d'un polígon semblant amb raó $k=2{,}5$.\recorda{Si la raó de semblança és $k$, els perímetres es multipliquen per $k$ i les àrees, per $k^2$.}

\ex{2} Un quadrat de 6 cm de costat es redueix amb raó $k=\frac13$. Calcula el costat i l'àrea del quadrat nou. Quantes vegades és més menuda l'àrea?

\ex{2} Un triangle té costats de 6, 8 i 10 cm. Un triangle semblant té 72 cm de perímetre. Quant fan els seus costats?

\ex{3} Dos quadrats tenen àrees de 9 cm² i 81 cm². Quant fa el costat de cadascun? Quina és la raó de semblança?

\ex{3} Ampliem una foto de 10 cm × 15 cm i la foto ampliada té 600 cm² d'àrea. Quina és la raó de semblança? Quines dimensions té la foto ampliada?

\ex{3} Un triangle rectangle té catets de 3 cm i 4 cm, i la hipotenusa fa 5 cm. Un triangle semblant té una hipotenusa de 15 cm. Calcula els seus catets i compara les àrees dels dos triangles.


% =====================================================
\seccio{2}{Plànols, mapes i maquetes}

\ex{1} Calcula la distància real i expressa-la en m o en km:\begin{tasks}(2)
  \task 7 cm en un mapa $1:50\,000$
  \task 4 cm en un plànol $1:100$
  \task 2,5 cm en un mapa $1:200\,000$
  \task 12 cm en un plànol $1:250$
\end{tasks}\recorda{Escala $1:n$: 1 cm del plànol són $n$ cm reals. \ Real $=$ plànol $\times\,n$; \ plànol $=$ real $:\,n$. \ 1 m $=100$ cm; \ 1 km $=100\,000$ cm.}

\ex{1} Quina longitud tindrà en el plànol o en el mapa?\begin{tasks}(1)
  \task Una paret de 8 m, en un plànol $1:200$.
  \task Una carretera de 15 km, en un mapa $1:500\,000$.
  \task Un carrer de 120 m, en un plànol $1:2000$.
\end{tasks}

\ex{1} Calcula l'escala:\begin{tasks}(1)
  \task 3 cm del mapa representen 1,5 km.
  \task 5 cm del plànol representen 10 m.
  \task 2 cm del mapa representen 40 km.
\end{tasks}

\ex{2} Aquesta és l'escala gràfica d'un mapa (cada tram fa 1 cm). Quina és l'escala numèrica? Quina distància real hi ha entre dos punts que estan a 6,5 cm en el mapa?\par\smallskip\begin{center}\begin{tikzpicture}[x=1cm,y=1cm,baseline=(current bounding box.center),line join=round]
\draw[fill=teal!12,draw=teal,thick] (0.000,0.000) -- (4.000,0.000) -- (4.000,0.300) -- (0.000,0.300) -- cycle;
\draw[fill=accent!22,draw=accent,thick] (0.000,0.000) -- (1.000,0.000) -- (1.000,0.300) -- (0.000,0.300) -- cycle;
\draw[fill=accent!22,draw=accent,thick] (2.000,0.000) -- (3.000,0.000) -- (3.000,0.300) -- (2.000,0.300) -- cycle;
\node[font=\small,inner sep=1pt] at (0.000,-0.350) {0};
\node[font=\small,inner sep=1pt] at (1.000,-0.350) {0,5};
\node[font=\small,inner sep=1pt] at (2.000,-0.350) {1};
\node[font=\small,inner sep=1pt] at (3.000,-0.350) {1,5};
\node[font=\small,inner sep=1pt] at (4.250,-0.350) {2 km};
\end{tikzpicture}\end{center}

\ex{2} En el plànol d'una casa a escala $1:50$, un dormitori fa 8 cm × 6,4 cm. Quines dimensions reals té? Quina superfície?

\ex{2} En un mapa a escala $1:25\,000$, el camí d'una excursió fa 18 cm. Quants km té l'excursió? Quant de temps tardarem si caminem a 5 km/h?

\ex{2} Una maqueta d'un avió està feta a escala $1:72$. L'avió real fa 36 m de llarg i 28,8 m d'ala a ala. Quines mesures té la maqueta?

\ex{2} En una maqueta a escala $1:500$, un edifici fa 9 cm d'altura. Quina altura té l'edifici real?

\ex{3} Dues ciutats que estan a 180 km es troben a 6 cm en un mapa. Quina és l'escala del mapa? A quina distància real està una altra ciutat que en el mapa és a 2,5 cm?

\ex{3} En un plànol a escala $1:100$, una parcel·la rectangular fa 25 cm × 16 cm. Quina superfície real té? Si el metre quadrat costa 150~€, quant costa la parcel·la?

\ex{3} En un llibre de ciències hi ha el dibuix d'una formiga a escala $20:1$ (és a dir, ampliat 20 vegades). Si en el dibuix la formiga fa 8 cm, quant fa la formiga de veritat?


% =====================================================
\seccio{3}{El teorema de Tales}

\ex{1} Calcula $x$ en cada proporció:\begin{tasks}(4)
  \task $\dfrac34=\dfrac{x}{12}$
  \task $\dfrac5x=\dfrac{10}{14}$
  \task $\dfrac x6=\dfrac{8}{12}$
  \task $\dfrac{2{,}5}{5}=\dfrac3x$
\end{tasks}

\ex{1} Tres rectes paral·leles tallen dues rectes $a$ i $b$. Sobre $a$ determinen $\overline{AB}=2$ cm i $\overline{BC}=5$ cm, i sobre $b$, $\overline{A'B'}=3$ cm. Calcula $\overline{B'C'}$.\recorda{\textbf{Teorema de Tales:} si unes rectes paral·leles tallen dues rectes, els segments que determinen sobre l'una són proporcionals als de l'altra: $\frac{\overline{AB}}{\overline{BC}}=\frac{\overline{A'B'}}{\overline{B'C'}}$.}

\ex{1} Les rectes \textcolor{navy}{blau fosc} són paral·leles. Calcula $x$.\par\smallskip\begin{center}\begin{tikzpicture}[x=0.55cm,y=0.55cm,baseline=(current bounding box.center),line join=round]
\draw[draw=navy,thick] (-1.200,0.000) -- (5.503,0.000);
\draw[draw=navy,thick] (-1.200,3.000) -- (5.503,3.000);
\draw[draw=navy,thick] (-1.200,5.000) -- (5.503,5.000);
\draw[thick] (-0.490,-0.700) -- (3.991,5.700);
\draw[thick] (3.502,-0.700) -- (4.401,5.700);
\draw[draw=accent,line width=1.6pt] (4.022,3.000) -- (4.303,5.000);
\fill (0.000,0.000) circle (1.5pt);
\fill (2.101,3.000) circle (1.5pt);
\fill (3.501,5.000) circle (1.5pt);
\fill (3.600,0.000) circle (1.5pt);
\fill (4.022,3.000) circle (1.5pt);
\fill (4.303,5.000) circle (1.5pt);
\node[font=\small,inner sep=2pt,above left] at (0.000,0.000) {$A$};
\node[font=\small,inner sep=2pt,above left] at (2.101,3.000) {$B$};
\node[font=\small,inner sep=2pt,above left] at (3.501,5.000) {$C$};
\node[font=\small,inner sep=2pt,above right] at (3.600,0.000) {$A'$};
\node[font=\small,inner sep=2pt,above right] at (4.022,3.000) {$B'$};
\node[font=\small,inner sep=2pt,above right] at (4.303,5.000) {$C'$};
\node[font=\small,inner sep=2pt,left] at (1.050,1.500) {$6$};
\node[font=\small,inner sep=2pt,left] at (2.801,4.000) {$4$};
\node[font=\small,inner sep=2pt,right] at (3.811,1.500) {$4{,}5$};
\node[font=\small,inner sep=2pt,right,text=accent] at (4.162,4.000) {$x$};
\end{tikzpicture}\end{center}

\ex{2} Les rectes \textcolor{navy}{blau fosc} són paral·leles. Calcula $x$.\par\smallskip\begin{center}\begin{tikzpicture}[x=0.55cm,y=0.55cm,baseline=(current bounding box.center),line join=round]
\draw[draw=navy,thick] (-1.200,0.000) -- (5.432,0.000);
\draw[draw=navy,thick] (-1.200,1.500) -- (5.432,1.500);
\draw[draw=navy,thick] (-1.200,4.500) -- (5.432,4.500);
\draw[thick] (-0.490,-0.700) -- (3.641,5.200);
\draw[thick] (3.502,-0.700) -- (4.331,5.200);
\draw[draw=accent,line width=1.6pt] (0.000,0.000) -- (1.050,1.500);
\fill (0.000,0.000) circle (1.5pt);
\fill (1.050,1.500) circle (1.5pt);
\fill (3.151,4.500) circle (1.5pt);
\fill (3.600,0.000) circle (1.5pt);
\fill (3.811,1.500) circle (1.5pt);
\fill (4.232,4.500) circle (1.5pt);
\node[font=\small,inner sep=2pt,above left] at (0.000,0.000) {$A$};
\node[font=\small,inner sep=2pt,above left] at (1.050,1.500) {$B$};
\node[font=\small,inner sep=2pt,above left] at (3.151,4.500) {$C$};
\node[font=\small,inner sep=2pt,above right] at (3.600,0.000) {$A'$};
\node[font=\small,inner sep=2pt,above right] at (3.811,1.500) {$B'$};
\node[font=\small,inner sep=2pt,above right] at (4.232,4.500) {$C'$};
\node[font=\small,inner sep=2pt,left,text=accent] at (0.525,0.750) {$x$};
\node[font=\small,inner sep=2pt,left] at (2.101,3.000) {$8$};
\node[font=\small,inner sep=2pt,right] at (3.705,0.750) {$3$};
\node[font=\small,inner sep=2pt,right] at (4.022,3.000) {$6$};
\end{tikzpicture}\end{center}

\ex{2} Amb rectes paral·leles, $\overline{AB}=3{,}6$, $\overline{BC}=4{,}8$ i $\overline{B'C'}=6$. Calcula $\overline{A'B'}$ i $\overline{A'C'}$.

\ex{2} Unes rectes tallen dues secants i hi determinen aquests segments. Poden ser paral·leles les rectes?\begin{tasks}(2)
  \task 3 i 4,5 en una; 4 i 6 en l'altra
  \task 2 i 5 en una; 3 i 7 en l'altra
\end{tasks}

\ex{2} Quatre rectes paral·leles determinen sobre una recta $\overline{AB}=2$, $\overline{BC}=3$ i $\overline{CD}=4$, i sobre una altra, $\overline{A'B'}=3$. Calcula $\overline{B'C'}$ i $\overline{C'D'}$.

\ex{2} Dividim un segment de 21 cm: \ a) en 3 parts iguals; \ b) en parts proporcionals a 1, 2 i 4. Quant fa cada part?

\ex{3} Tres solars tenen les partions paral·leles i donen a dos carrers. Al carrer Major, les façanes fan 18 m, 24 m i 30 m. Al carrer Nou, les tres façanes fan en total 84 m. Quant fa cada façana al carrer Nou?

\ex{3} Amb rectes paral·leles, $\overline{AB}=x$, $\overline{BC}=x+2$, $\overline{A'B'}=6$ i $\overline{B'C'}=9$. Calcula $x$.

\ex{3} Explica com dividiries un segment de 10 cm en 3 parts iguals amb regle i escaire, sense mesurar. Quant farà cada part?


% =====================================================
\seccio{4}{Aplicació a la semblança de triangles}

\ex{1} Són semblants els triangles que tenen aquests costats? Si ho són, calcula la raó de semblança.\begin{tasks}(1)
  \task 4, 5, 6 \ i \ 8, 10, 12
  \task 3, 3, 4 \ i \ 6, 6, 9
  \task 5, 12, 13 \ i \ 7,5, 18, 19,5
\end{tasks}\recorda{Dos triangles són semblants si tenen \textbf{dos angles iguals}, o els \textbf{tres costats proporcionals}, o \textbf{dos costats proporcionals i l'angle que formen igual}. Per a comparar costats, ordena'ls de menor a major.}

\ex{1} Són semblants aquests triangles?\begin{tasks}(1)
  \task Un té angles de 45° i 60°, i l'altre, de 60° i 75°.
  \task Un té angles de 50° i 50°, i l'altre, de 50° i 90°.
\end{tasks}

\ex{1} En la figura, $B'C'\parallel BC$. Calcula $x$.\par\smallskip\begin{center}\begin{tikzpicture}[x=0.5cm,y=0.5cm,baseline=(current bounding box.center),line join=round]
\draw[fill=teal!12,draw=teal,thick] (0.000,0.000) -- (-2.817,-5.298) -- (2.644,-5.938) -- cycle;
\draw[fill=accent!22,draw=accent,thick] (0.000,0.000) -- (-1.878,-3.532) -- (1.763,-3.959) -- cycle;
\draw[draw=accent,line width=1.6pt] (-1.878,-3.532) -- (1.763,-3.959);
\draw[draw=teal,line width=1.6pt] (-2.817,-5.298) -- (2.644,-5.938);
\fill (0.000,0.000) circle (1.5pt);
\fill (-2.817,-5.298) circle (1.5pt);
\fill (2.644,-5.938) circle (1.5pt);
\fill (-1.878,-3.532) circle (1.5pt);
\fill (1.763,-3.959) circle (1.5pt);
\node[font=\small,inner sep=2pt,above] at (0.000,0.000) {$A$};
\node[font=\small,inner sep=2pt,left] at (-2.817,-5.298) {$B$};
\node[font=\small,inner sep=2pt,right] at (2.644,-5.938) {$C$};
\node[font=\small,inner sep=2pt,left] at (-1.878,-3.532) {$B'$};
\node[font=\small,inner sep=2pt,right] at (1.763,-3.959) {$C'$};
\node[font=\small,inner sep=2pt,left] at (-0.939,-1.766) {$4$};
\node[font=\small,inner sep=2pt,left] at (-2.347,-4.415) {$2$};
\node[font=\small,inner sep=2pt,below] at (-0.087,-5.618) {$9$};
\node[font=\small,inner sep=2pt,above,text=accent] at (-0.058,-3.745) {$x$};
\end{tikzpicture}\end{center}

\ex{2} En la figura, $B'C'\parallel BC$. Calcula $x$.\par\smallskip\begin{center}\begin{tikzpicture}[x=0.5cm,y=0.5cm,baseline=(current bounding box.center),line join=round]
\draw[fill=teal!12,draw=teal,thick] (0.000,0.000) -- (-2.817,-5.298) -- (2.644,-5.938) -- cycle;
\draw[fill=accent!22,draw=accent,thick] (0.000,0.000) -- (-1.878,-3.532) -- (1.763,-3.959) -- cycle;
\draw[draw=accent,line width=1.6pt] (-1.878,-3.532) -- (1.763,-3.959);
\draw[draw=teal,line width=1.6pt] (-2.817,-5.298) -- (2.644,-5.938);
\fill (0.000,0.000) circle (1.5pt);
\fill (-2.817,-5.298) circle (1.5pt);
\fill (2.644,-5.938) circle (1.5pt);
\fill (-1.878,-3.532) circle (1.5pt);
\fill (1.763,-3.959) circle (1.5pt);
\node[font=\small,inner sep=2pt,above] at (0.000,0.000) {$A$};
\node[font=\small,inner sep=2pt,left] at (-2.817,-5.298) {$B$};
\node[font=\small,inner sep=2pt,right] at (2.644,-5.938) {$C$};
\node[font=\small,inner sep=2pt,left] at (-1.878,-3.532) {$B'$};
\node[font=\small,inner sep=2pt,right] at (1.763,-3.959) {$C'$};
\node[font=\small,inner sep=2pt,left] at (-0.939,-1.766) {$6$};
\node[font=\small,inner sep=2pt,left] at (-2.347,-4.415) {$3$};
\node[font=\small,inner sep=2pt,right] at (0.881,-1.979) {$8$};
\node[font=\small,inner sep=2pt,right,text=accent] at (2.203,-4.948) {$x$};
\end{tikzpicture}\end{center}

\ex{2} En la figura, $B'C'\parallel BC$. Calcula $x$.\par\smallskip\begin{center}\begin{tikzpicture}[x=0.5cm,y=0.5cm,baseline=(current bounding box.center),line join=round]
\draw[fill=teal!12,draw=teal,thick] (0.000,0.000) -- (-2.817,-5.298) -- (2.644,-5.938) -- cycle;
\draw[fill=accent!22,draw=accent,thick] (0.000,0.000) -- (-1.408,-2.649) -- (1.322,-2.969) -- cycle;
\draw[draw=accent,line width=1.6pt] (-1.408,-2.649) -- (1.322,-2.969);
\draw[draw=teal,line width=1.6pt] (-2.817,-5.298) -- (2.644,-5.938);
\fill (0.000,0.000) circle (1.5pt);
\fill (-2.817,-5.298) circle (1.5pt);
\fill (2.644,-5.938) circle (1.5pt);
\fill (-1.408,-2.649) circle (1.5pt);
\fill (1.322,-2.969) circle (1.5pt);
\node[font=\small,inner sep=2pt,above] at (0.000,0.000) {$A$};
\node[font=\small,inner sep=2pt,left] at (-2.817,-5.298) {$B$};
\node[font=\small,inner sep=2pt,right] at (2.644,-5.938) {$C$};
\node[font=\small,inner sep=2pt,left] at (-1.408,-2.649) {$B'$};
\node[font=\small,inner sep=2pt,right] at (1.322,-2.969) {$C'$};
\node[font=\small,inner sep=2pt,left] at (-0.704,-1.324) {$4$};
\node[font=\small,inner sep=2pt,left] at (-2.113,-3.973) {$4$};
\node[font=\small,inner sep=2pt,above] at (-0.043,-2.809) {$5$};
\node[font=\small,inner sep=2pt,below,text=accent] at (-0.087,-5.618) {$x$};
\end{tikzpicture}\end{center}

\ex{2} Un triangle rectangle té un angle de 40°. Un altre triangle rectangle té un angle de 50°. Són semblants? Per què?

\ex{2} Els costats d'un triangle fan 6, 8 i 9 cm. En un triangle semblant, el costat menut fa 9 cm. Quant fan els altres dos costats?

\ex{2} En els triangles $ABC$ i $A'B'C'$, $\hat A=\hat A'=60^\circ$, $\overline{AB}=4$, $\overline{AC}=6$, $\overline{A'B'}=6$ i $\overline{A'C'}=9$. Són semblants? Si $\overline{BC}=5{,}3$, quant fa $\overline{B'C'}$?

\ex{3} Els costats d'un triangle fan 4, 6 i 8 cm. Un triangle semblant té 45 cm de perímetre. Quant fan els seus costats?

\ex{3} En un triangle $ABC$, tracem $B'C'\parallel BC$ amb $\overline{AB'}=5$, $\overline{B'B}=10$, $\overline{BC}=12$ i $\overline{AC'}=6$. Calcula $\overline{B'C'}$, $\overline{C'C}$ i el perímetre del trapezi $B'BCC'$.

\ex{3} Un triangle té una base de 4 cm i una altura de 3 cm. En un triangle semblant, la base fa 10 cm. Calcula'n l'altura i l'àrea. Per quant s'ha multiplicat l'àrea?


% =====================================================
\seccio{5}{Problemes amb semblança de triangles}

\ex{1} Un arbre fa una ombra de 9 m. Al mateix moment, un pal de 2 m fa una ombra d'1,5 m. Quina altura té l'arbre?\recorda{Fes un dibuix, busca dos triangles semblants (amb els mateixos angles o en posició de Tales) i escriu la proporció entre costats homòlegs. Al mateix moment, els rajos del Sol són paral·lels: altura i ombra són proporcionals.}

\ex{1} Un edifici de 21 m fa una ombra de 14 m. Quina ombra fa, al mateix moment, una persona d'1,5 m?

\ex{1} Calcula l'altura $H$ de la torre.\par\smallskip\begin{center}\begin{tikzpicture}[x=0.42cm,y=0.42cm,baseline=(current bounding box.center),line join=round]
\draw[muted] (-0.600,0.000) -- (8.400,0.000);
\draw[fill=accent!22,draw=accent,thick] (0.000,0.000) -- (0.000,1.200) -- (0.900,0.000) -- cycle;
\draw[thin] (0.524,0.000) -- (0.524,0.524) -- (0.000,0.524);
\draw[fill=teal!12,draw=teal,thick] (3.100,0.000) -- (3.100,6.000) -- (7.600,0.000) -- cycle;
\draw[thin] (3.624,0.000) -- (3.624,0.524) -- (3.100,0.524);
\draw[line width=2pt,draw=muted] (0.000,0.000) -- (0.000,1.200);
\draw[line width=2pt,draw=muted] (3.100,0.000) -- (3.100,6.000);
\draw[dashed,draw=muted] (0.000,1.200) -- (0.900,0.000);
\draw[dashed,draw=muted] (3.100,6.000) -- (7.600,0.000);
\node[font=\small,inner sep=2pt,left] at (0.000,0.600) {$1{,}2$\,m};
\node[font=\small,inner sep=2pt,below] at (0.450,0.000) {$0{,}9$\,m};
\node[font=\small,inner sep=2pt,left,text=accent] at (3.100,3.000) {$H$};
\node[font=\small,inner sep=2pt,below] at (5.350,0.000) {$15$\,m};
\end{tikzpicture}\end{center}

\ex{2} La Clara té els ulls a 1,6 m de terra. Posa un mirall a terra, a 2 m d'ella i a 15 m d'un edifici, i hi veu just el capdamunt de l'edifici. Quina altura té l'edifici?\par\smallskip\begin{center}\begin{tikzpicture}[x=0.22cm,y=0.22cm,baseline=(current bounding box.center),line join=round]
\draw[muted] (-0.500,0.000) -- (17.500,0.000);
\draw[line width=2pt,draw=muted] (0.000,0.000) -- (0.000,1.600);
\draw[line width=2pt,draw=muted] (17.000,0.000) -- (17.000,12.000);
\draw[dashed,draw=muted] (0.000,1.600) -- (2.000,0.000);
\draw[dashed,draw=muted] (2.000,0.000) -- (17.000,12.000);
\draw[fill=accent!22,draw=accent,thick] (0.000,0.000) -- (0.000,1.600) -- (2.000,0.000) -- cycle;
\draw[fill=teal!12,draw=teal,thick] (2.000,0.000) -- (17.000,0.000) -- (17.000,12.000) -- cycle;
\draw[thin] (1.000,0.000) -- (1.000,1.000) -- (0.000,1.000);
\draw[thin] (16.000,0.000) -- (16.000,1.000) -- (17.000,1.000);
\fill (2.000,0.000) circle (1.5pt);
\node[font=\small,inner sep=2pt,left] at (0.000,0.800) {$1{,}6$\,m};
\node[font=\small,inner sep=2pt,below] at (1.000,0.000) {$2$\,m};
\node[font=\small,inner sep=2pt,below] at (9.500,0.000) {$15$\,m};
\node[font=\small,inner sep=2pt,right,text=accent] at (17.000,6.000) {$H$};
\end{tikzpicture}\end{center}

\ex{2} Una rampa puja 0,9 m d'altura en 6 m horitzontals. A quina altura està un punt de la rampa que és a 4 m (horitzontals) del començament?

\ex{2} Una persona d'1,5 m està a 2 m d'un fanal de 4,5 m d'altura. Quant fa la seua ombra?

\ex{2} Una escala recolzada en una paret arriba a 4 m d'altura, i el seu peu està a 1,5 m de la paret. A quina altura $h$ està el punt de l'escala que és a 0,6 m de la paret?\par\smallskip\begin{center}\begin{tikzpicture}[x=0.75cm,y=0.75cm,baseline=(current bounding box.center),line join=round]
\draw[muted] (-0.400,0.000) -- (1.900,0.000);
\draw[line width=2pt,draw=muted] (0.000,0.000) -- (0.000,4.400);
\draw[fill=teal!12,draw=teal,thick] (0.000,0.000) -- (1.500,0.000) -- (0.000,4.000) -- cycle;
\draw[dashed,draw=muted] (0.600,0.000) -- (0.600,2.400);
\fill (0.600,2.400) circle (1.5pt);
\draw[thin] (0.293,0.000) -- (0.293,0.293) -- (0.000,0.293);
\node[font=\small,inner sep=2pt,left] at (0.000,2.000) {$4$\,m};
\node[font=\small,inner sep=2pt,right,text=accent] at (0.600,1.200) {$h$};
\draw[dashed,draw=muted] (0.000,-0.350) -- (0.600,-0.350);
\node[font=\small,inner sep=1pt] at (0.300,-0.600) {0,6 m};
\draw[dashed,draw=muted] (0.000,-0.950) -- (1.500,-0.950);
\node[font=\small,inner sep=1pt] at (0.750,-1.200) {1,5 m};
\end{tikzpicture}\end{center}

\ex{2} Amb el braç estirat, tapem exactament un arbre de 6 m d'altura amb un regle vertical de 10 cm que està a 50 cm dels ulls. A quina distància està l'arbre?

\ex{3} Per a mesurar l'amplària $x$ d'un riu s'han pres les mesures de la figura ($AB\parallel DE$ i $A$, $C$, $E$ alineats). Quant fa d'ample el riu?\par\smallskip\begin{center}\begin{tikzpicture}[x=0.3cm,y=0.3cm,baseline=(current bounding box.center),line join=round]
\draw[fill=teal!14,draw=none] (-1.000,0.300) -- (12.000,0.300) -- (12.000,7.200) -- (-1.000,7.200) -- cycle;
\draw[fill=teal!12,draw=teal,thick] (0.000,0.000) -- (0.000,7.500) -- (9.000,0.000) -- cycle;
\draw[fill=accent!22,draw=accent,thick] (9.000,0.000) -- (12.000,0.000) -- (12.000,-2.500) -- cycle;
\draw[muted] (-1.000,0.000) -- (12.500,0.000);
\draw[thin] (0.733,0.000) -- (0.733,0.733) -- (0.000,0.733);
\draw[thin] (11.267,0.000) -- (11.267,-0.733) -- (12.000,-0.733);
\node[font=\small,inner sep=2pt,left] at (0.000,7.500) {$A$};
\node[font=\small,inner sep=2pt,left] at (0.000,0.000) {$B$};
\node[font=\small,inner sep=2pt,above left] at (9.000,0.000) {$C$};
\node[font=\small,inner sep=2pt,right] at (12.000,0.000) {$D$};
\node[font=\small,inner sep=2pt,right] at (12.000,-2.500) {$E$};
\node[font=\small,inner sep=2pt,below] at (4.500,0.000) {$9$\,m};
\node[font=\small,inner sep=2pt,above] at (10.500,0.000) {$3$\,m};
\node[font=\small,inner sep=2pt,right] at (12.000,-1.250) {$2{,}5$\,m};
\node[font=\small,inner sep=2pt,left,text=accent] at (0.000,3.750) {$x$};
\end{tikzpicture}\end{center}

\ex{3} Tales, que feia 1,8 m d'altura, va mesurar la seua ombra: 2,7 m. Al mateix moment, l'ombra de la piràmide, comptada des del centre de la base, feia 219 m. Quina altura tenia la piràmide?

\ex{3} Un nen d'1,2 m es col·loca de manera que la punta de la seua ombra coincideix amb la de l'ombra d'un arbre. L'ombra del nen fa 1,6 m i el nen està a 6,4 m de l'arbre. Quina altura té l'arbre?

\vfill
\begin{center}\color{muted}\small\faCheckCircle[regular]\ \ Ànim! Comprova els teus resultats amb les eines interactives dels apunts de la unitat.\end{center}

\end{document}
